% ==================================================================
% id:       func-011
% topic:    מספרים ממשיים
% subtopic: הוכחת אי־רציונליות
% solution: full
% ==================================================================
{\bf
שאלת אתגר
}
\begin{enumerate}
\item
הוכיחו כי
$\sqrt{3} \notin \mathbb{Q}$.
\hint{חקו את ההוכחה ש-$\sqrt2\notin\mathbb{Q}$.}
\begin{solbox}
ההוכחה תחקה את ההוכחה ש-
$\sqrt2\notin\mathbb{Q}$.
נניח בשלילה כי
\begin{equation}\label{eq:sqrt3}
\sqrt3=\frac{m}{n}
\end{equation}
עבור
$m,n$
שלמים זרים (כלומר, ללא מחלק משותף פרט ל-$\pm1$)
ו-$n\neq0$.
נעלה את השוויון
\eqref{eq:sqrt3}
בריבוע ונקבל
$3=\frac{m^2}{n^2}$,
כלומר
$$.3n^2=m^2$$
אנו רואים כי
$m^2$
מתחלק ב-$3$
ולכן
$m$
מתחלק ב-$3$
(אחרת,
$m^2$
לא היה מתחלק ב-$3$).
לכן ניתן לכתוב
$m=3k$
עבור
$k\in\mathbb{Z}$.
נציב ונקבל
$3n^2=9k^2$,
כלומר,
$n^2=3k^2$.
כמקודם, נסיק כי
$n$
מתחלק ב-$3$.
זאת בסתירה להנחתנו כי
$m,n$
הם זרים.
\end{solbox}
\item
הוכיחו כי
$\sqrt{2} + \sqrt{3} \notin \mathbb{Q}$.

הדרכה:
חשבו את
$(\sqrt{2} + \sqrt{3})(\sqrt{2} - \sqrt{3}) $,
והראו כי אם
$a, b \in \rr$
מקיימים
$ab = -1$
אז או ששניהם רציונליים או ששניהם אי רציונליים. לבסוף, הסבירו למה לא ייתכן ש
$\sqrt{2} + \sqrt{3} , \sqrt{2} - \sqrt{3} $
שניהם רציונליים.
\hint{עקבו אחרי ההדרכה. כדי להראות ש-$\sqrt2+\sqrt3$ ו-$\sqrt2-\sqrt3$ אינם שניהם רציונליים, הסתכלו על הסכום שלהם.}
\begin{solbox}
נראה כי
$\sqrt2+\sqrt3\notin\mathbb{Q}$
בשתי דרכים שונות:
\begin{itemize}
\item
דרך א': ניעזר בטריק של הכפלה בצמוד! נסמן
$$.a=\sqrt2+\sqrt3\ ,\ b=\sqrt2-\sqrt3$$
ונשים לב כי
$a\cdot b=-1$
לכן
$b=-\frac1a$.
מכאן ניתן להסיק כי אילו
$a\in\mathbb{Q}$
אז גם
$b\in\mathbb{Q}$
ואז ודאי
$\sqrt2=\frac12(a+b)\in\mathbb{Q}$,
וזוהי כמובן סתירה.
\item
ניתן להוכיח את העובדה גם בשלילה, בדומה לסעיף הקודם. כלומר, מניחים כי
$$.\sqrt2+\sqrt3=\frac{m}{n},\ m,n\in\mathbb{Z},\ n\neq0$$
מעלים בריבוע ומבטאים את
$\sqrt6$
כמנה של מספרים שלמים. מכאן די להוכיח כי
$\sqrt6\notin\mathbb{Q}$,
בצורה דומה לפתרון של הסעיף הקודם.
\end{itemize}
\end{solbox}
\end{enumerate}
